\documentclass[10pt,a4paper]{article}
\usepackage[utf8]{inputenc}
\usepackage[T1]{fontenc}
\usepackage[french]{babel}
\usepackage{geometry,amsmath,amssymb,array,multicol,fancyhdr,graphicx,tikz}
\geometry{top=1.25cm,bottom=1.25cm,left=1.2cm,right=1.2cm}
\setlength{\parindent}{0pt}
\setlength{\parskip}{2pt}
\setlength{\columnsep}{0.65cm}
\setlength{\columnseprule}{0.25pt}
\raggedcolumns
\setlength{\emergencystretch}{2em}
\pagestyle{fancy}\fancyhf{}
\lhead{\small Mathématiques - 4e}
\rhead{\small Support long — produits remarquables}
\cfoot{\thepage}
\renewcommand{\headrulewidth}{0.3pt}
\renewcommand{\footrulewidth}{0.3pt}
\newcommand{\SequenceTitle}[1]{\par\medskip\hrule\smallskip\textbf{\large #1}\smallskip\hrule\medskip}
\newcommand{\ActivityTitle}[2]{\par\medskip\textbf{Activité #1 : #2}\par\nobreak\smallskip}
\newcommand{\InstructionTitle}[2]{\par\smallskip\textbf{Consigne #1#2}\par\nobreak}
\newcommand{\PedagogicalBlock}[3]{\par\smallskip\textbf{#1#2}\par\nobreak #3\par}
\newcommand{\Separator}{\par\smallskip\noindent\dotfill\par\smallskip}
% Keep ordinary cutouts intact; very long authored blocks remain breakable.
\newcommand{\PrintableBlock}[1]{\par\setbox0=\vbox{\hsize=\columnwidth\linewidth=\columnwidth #1}\ifdim\dimexpr\ht0+\dp0\relax<0.85\textheight\box0\else #1\fi\par}
\begin{document}
\begin{center}\textbf{\Large Support long — produits remarquables}\\
\small SA1 Configurations du plan - Séquence 8 Calculs sur les expressions algébriques
\end{center}
\begin{multicols}{2}

\ActivityTitle{1}{Découverte — parcelle de Fofo}
\PrintableBlock{
\par\smallskip\textbf{Situation de départ}\par Fofo dispose d'une parcelle rectangulaire dont les côtés sont partagés en longueurs $x$, $z$ d'une part et $y$, $b$ d'autre part. Exprimer son aire de deux façons.\par
\Separator
}
\PrintableBlock{
\InstructionTitle{1}{ :}Calcule l'aire du grand rectangle, puis la somme des aires des quatre rectangles. Compare les résultats et complète : $(x+z)(y+b)=\ldots$
\Separator
}
\ActivityTitle{2}{Établir les produits remarquables}
\PrintableBlock{
\InstructionTitle{2}{ :}Développe et réduis $(a+b)^2$, $(a-b)^2$ et $(a+b)(a-b)$, puis formule les égalités obtenues.
\Separator
}
\PrintableBlock{
\PedagogicalBlock{Propriété}{ :}{\emph{Pour tous nombres $a$ et $b$ : $(a+b)^2=a^2+2ab+b^2$ ; $(a-b)^2=a^2-2ab+b^2$ ; $(a+b)(a-b)=a^2-b^2$.}}
\Separator
}
\PrintableBlock{
\PedagogicalBlock{Application}{ :}{Développe : $A=(2x+1)^2$, $B=(3x-2)^2$, $C=(4x-3)(4x+3)$. Factorise : $D=x^2+6x+9$, $E=4x^2-9$, $F=4x^2-4x+1$.}
\Separator
}
\PrintableBlock{
\InstructionTitle{3}{ : Recherche guidée}Pour $x=2$ et $y=3$, calcule séparément $(x+y)^2$ puis $x^2+2xy+y^2$. Recommence avec $x=5$ et $y=-1$. Compare chaque paire de résultats et explique, en une phrase complète, ce que montrent ces vérifications numériques. Présente tous les calculs intermédiaires et contrôle le signe du double produit.
\Separator
}
\PrintableBlock{
\InstructionTitle{4}{ : Développement raisonné}Développe et réduis les expressions $G=(3x+4)^2$, $H=(5-2x)^2$ et $I=(7x-1)(7x+1)$. Pour chaque expression, indique l'identité remarquable utilisée, identifie les deux termes, écris l'étape de substitution puis donne la forme réduite. Vérifie enfin un des résultats pour $x=2$.
\Separator
}
\PrintableBlock{
\InstructionTitle{5}{ : Factorisation raisonnée}Factorise $J=x^2+10x+25$, $K=9x^2-16$ et $L=4x^2-12x+9$. Avant de calculer, repère les carrés et le double produit éventuel. Écris l'identité remarquable dans le sens utilisé, justifie ton choix puis vérifie chaque factorisation en redéveloppant le produit obtenu.
\Separator
}
\PrintableBlock{
\InstructionTitle{6}{ : Choisir une stratégie}Sans effectuer immédiatement tous les calculs, indique pour chaque expression s'il est plus efficace de développer, de factoriser ou de conserver la forme donnée : $(x+8)^2$, $x^2-64$, $4x^2+12x+9$, $(2x-3)(2x+3)$. Justifie chaque choix, puis effectue les transformations demandées et compare les nombres d'opérations.
\Separator
}
\PrintableBlock{
\InstructionTitle{7}{ : Problème de synthèse}Un jardin carré de côté $(x+3)$ mètres contient une allée carrée de côté $x$ mètres. Exprime l'aire restante sous une forme développée puis sous une forme factorisée. Interprète géométriquement chaque terme, calcule cette aire pour $x=7$ et rédige une conclusion avec l'unité correcte.
\Separator
}
\PrintableBlock{
\PedagogicalBlock{Définition}{ : Identité remarquable}{\emph{Une identité remarquable est une égalité vraie pour toutes les valeurs admissibles des lettres. Elle permet de développer ou de factoriser rapidement certaines expressions algébriques.}}
\Separator
}
\PrintableBlock{
\PedagogicalBlock{Méthode}{ : Reconnaître puis transformer}{1. Repérer les carrés des deux termes. 2. Chercher le double produit $2ab$ et contrôler son signe. 3. Choisir l'identité correspondante. 4. Remplacer $a$ et $b$ sans oublier les parenthèses. 5. Réduire et vérifier le résultat en redéveloppant ou avec une valeur numérique simple.}
\Separator
}
\PrintableBlock{
\PedagogicalBlock{Remarque}{ : Attention au signe}{Dans $(a-b)^2$, le terme du milieu est $-2ab$, mais le dernier terme reste $+b^2$. Le carré d'un nombre négatif est positif.}
\Separator
}
\PrintableBlock{
\par\textbf{Schéma de la parcelle}\par\begin{center}\resizebox{\columnwidth}{!}{\begin{tikzpicture}[scale=0.75]\draw[thick] (0,0) rectangle (8,4);\draw (5,0)--(5,4);\draw (0,2.5)--(8,2.5);\node at (2.5,1.25) {$xy$};\node at (6.5,1.25) {$zy$};\node at (2.5,3.25) {$xb$};\node at (6.5,3.25) {$zb$};\draw[<->] (0,-0.4)--node[below] {$x$}(5,-0.4);\draw[<->] (5,-0.4)--node[below] {$z$}(8,-0.4);\end{tikzpicture}}\end{center}
\Separator
}
\end{multicols}
\end{document}